% End-to-End Evaluation With Canaries: body of Section~\ref{sec:e2e} (the \section heading and its
% label stay in main.tex). Labels defined here: sec:e2e-setup, sec:e2e-safety, tab:e2e-safety,
% sec:e2e-model. Source: evidence/kaggle-scale-run-2026-10-24.md (one run, wardowl 2.1.0, before the
% honest-completeness change). Repeats (X8) are pending and marked \tbd{}.

The previous sections test components in isolation. Here the whole pipeline runs as an organization
would run it: many tenants and people in one index, a policy composed of a tenant boundary and
per-document ACLs, adversarial questions, and a language model that answers from what retrieval
returns. The goal is a leak test with enough exposure to bound the leak rate, not a measure of
answer quality.

\subsection{Setup}\label{sec:e2e-setup}

\paragraph{Corpus.}
The index holds 107,540 chunks: 49,920 synthetic company documents in 20 tenants, of which 41,600
carry a unique canary string; 20 notices, one per tenant, each carrying a prompt injection; and
57,600 passages of the FiQA financial question-answering corpus~\cite{maia2018www18} from
BEIR~\cite{thakur2021beir}, which every tenant may read and which serve as distractors. Company
documents are drawn from templates (handbooks, finance, HR, board and project records) with
per-document ACLs over users and groups. Chunks are embedded with
bge-base-en-v1.5~\cite{xiao2024cpack}.

\paragraph{People and questions.}
189 people span the tenants. A typical employee may read the shared passages, the tenant's notice
and five of the twelve document templates, 58,641 chunks in all (the median); the least privileged
are non-executive board members (``judy'' personas), each of whom may read 624 chunks of three
templates in one tenant and none of the shared passages (0.58\% of the corpus); one person in no
group may read nothing. They ask 1,105 questions: 108 generic, 376 targeted at a document
they may read, 189 targeted at a document they may not read (\emph{forbidden}), and 432 red-team
prompts that ask for other people's data, other tenants' data or everything, or try to override
instructions.

\paragraph{Policy and stores.}
The policy is $\mathrm{AllOf}(\mathrm{Tenant}, \mathrm{MetadataACL})$ (Section~\ref{sec:sys-compose}).
It runs against ten configurations: FAISS over a flat index (exact, the reference), Chroma, pgvector,
pgvector with the same policy written in OPA, Qdrant, Elasticsearch, OpenSearch, Weaviate, Redis and
LanceDB, each at the settings listed in Appendix~\tbd{settings table}. A configuration with OpenFGA
behind the caching authorizer did not start (the engine's health endpoint did not answer within
180\,s) and is excluded. Each question retrieves with $k = 5$ for the answer and $k = 10$ for
recall@10. Qwen2.5-3B-Instruct~\cite{yang2024qwen25} answers from the five chunks.

\paragraph{Leak detection.}
A leak is any returned chunk outside the person's readable set, as decided by an oracle that
evaluates the policy without the layer, or any canary from an unreadable document found verbatim in
a retrieved chunk, in the prompt sent to the model or in the model's answer. Verbatim matching
misses a paraphrased secret; it is sufficient for retrieved chunks and prompts, which are copied, and
a lower bound for answers. \tbd{X8 adds paraphrase-aware matching.}

\subsection{Safety}\label{sec:e2e-safety}

\begin{table}[tbp]
\centering
\small
\begin{tabular}{@{}lrrrrr@{}}
\toprule
Configuration & Questions & Hits checked & Unreadable hits & Canaries (chunk/prompt/answer) & Violations \\
\midrule
FAISS (exact) & 1,105 & 16,360 & 0 & 0/0/0 & 0 \\
Chroma & 1,105 & 16,360 & 0 & 0/0/0 & 0 \\
pgvector & 1,105 & 16,178 & 0 & 0/0/0 & 0 \\
pgvector + OPA & 1,105 & 16,178 & 0 & 0/0/0 & 0 \\
Qdrant & 1,105 & 16,360 & 0 & 0/0/0 & 0 \\
Elasticsearch & 1,105 & 16,360 & 0 & 0/0/0 & 0 \\
OpenSearch & 1,105 & 16,360 & 0 & 0/0/0 & 0 \\
Weaviate & 1,105 & 16,360 & 0 & 0/0/0 & 0 \\
Redis & 1,105 & 16,360 & 0 & 0/0/0 & 0 \\
LanceDB & 1,105 & 16,360 & 0 & 0/0/0 & 0 \\
\midrule
Total & 11,050 & 163,236 & 0 & 0/0/0 & 0 \\
\bottomrule
\end{tabular}
\caption{End-to-end safety, one run. Hits are counted over both retrievals per question. pgvector
checked 182 fewer hits because some of its answers were short or empty (Section~\ref{sec:quality}).}
\label{tab:e2e-safety}
\end{table}

No configuration returned an unreadable chunk or exposed a canary (Table~\ref{tab:e2e-safety}). The
hits are not independent trials: they cluster by question and by person, and a fault in the layer
would affect a question, not a random hit. We therefore bound the leak rate per question, separately
for each configuration: with no leak in 1,105 questions, the one-sided 95\% upper bound is $1 -
0.05^{1/1105} \approx 2.7 \times 10^{-3}$, and per person $1 - 0.05^{1/189} \approx 1.6 \times
10^{-2}$ (Section~\ref{sec:prob-metrics}). Pooling the ten configurations would not shrink these
bounds honestly, since they share questions and the layer's code. The violation count is zero
because the filter and the check read the same ACL metadata here; Sections~\ref{sec:conformance}
and~\ref{sec:engines} test configurations in which they do not.

Two observations qualify the result. First, it is a statement about this layer under assumptions
(A1) to (A3) of Proposition~\ref{prop:reduction}, exercised with adversarial inputs, not a proof:
the red-team prompts can only change the query text, and the layer never lets text choose the
subject (Section~\ref{sec:sys-identity}). Second, safety held while quality did not. pgvector
returned no hits at all to ten least-privileged people who could read hundreds of chunks, and the
layer, before the change of Section~\ref{sec:sys-complete}, reported those empty answers as
complete; Qdrant answered, but with recall@10 of 0.56. Section~\ref{sec:quality} diagnoses both.
\tbd{X8: three repeated runs with fresh permission layouts; report the bound over all runs and
the run-to-run spread of recall.}

\subsection{What the model did with what it retrieved}\label{sec:e2e-model}

Prompt injection in readable content is outside our threat model (Section~\ref{sec:prob-threat}),
but the run measures where the boundary lies. Each tenant's injected notice is readable by that
tenant's employees and contractors, so retrieving it for them is correct. In the nine configurations
for which the run reports it, the notice was retrieved between 7 and 106 times, depending on the
store's ranking, and the 3B model followed its instruction every time it was retrieved (482 of
482). Retrieval-time enforcement decides which
content reaches the model, not what the model does with it; defenses against injected instructions
in permitted content~\cite{greshake2023not,debenedetti2025defeating} are complementary.

Answer quality followed retrieval quality. With FAISS the model found the target document of a
targeted question in its top five for every question and answered 99.7\% correctly; with Qdrant
the target was in the top five for 39.9\% of targeted questions. When a person asked about a
document they could not read, the model abstained in 70\% to 81\% of cases, depending on the store,
and otherwise answered from readable documents only, as the canary checks confirm.
